\documentclass[11pt]{article}
\usepackage[a4paper,margin=26mm]{geometry}
\usepackage{amsmath,amssymb,mathtools}
\usepackage{unicode-math}
\setmainfont{Latin Modern Roman}
\setmathfont{Latin Modern Math}
\usepackage{xurl,hyperref}
\hypersetup{hidelinks,pdftitle={Problem statement}}
\newcommand{\sref}[2]{\hyperlink{source:#1:#2}{[S#2]}}
\newcommand{\cataloguescope}[1]{#1}
\begin{document}
\section{Conjecture 2.5 — Vassiliev invariants distinguish oriented knots.}\label{statement}

\subsection{Definitions and mathematical statement}

\paragraph{Shared notation.}$\mathbb N=\{1,2,\ldots\}$, and $\mathbb Z,\mathbb Q,\mathbb R,\mathbb C$ denote the integers, rationals, reals and complex numbers. Any other conventions are specified locally.


A knot or link is a smooth oriented embedding of one or more circles in the oriented $3$-sphere, considered up to orientation-preserving ambient isotopy. Write $O$ for the unknot, which bounds an embedded disk. Let $R\mathcal K$ be the free $R$-module on oriented knot isotopy classes. Extend a knot invariant to singular knots with ordinary transverse double points by $v(K_\times)=v(K_+)-v(K_-)$. Let $F_d(R\mathcal K)$ be the submodule generated by the alternating resolution sums of knots with $d$ such points. An $R$-valued finite-type (Vassiliev) invariant of degree at most $d$ is an $R$-linear map $R\mathcal K\to R$ vanishing on $F_{d+1}(R\mathcal K)$. For coefficients in an abelian group $A$, use homomorphisms $\mathbb Z\mathcal K\to A$ that vanish on $F_{d+1}(\mathbb Z\mathcal K)$; this is the same finite-type definition with those coefficients.
\paragraph{Statement recorded in the supplied source.}
Do finite-type invariants distinguish oriented knots? Precisely, for any two nonisotopic oriented knots $K,K'$, is there a finite degree and a finite-type knot invariant $v$ with $v(K)\ne v(K')$? The degree can depend on the pair. The query concerns oriented isotopy classes, not their unoriented classes. \sref{statement}{2}
\subsection{Short English statement}
Can some finite-type invariant distinguish every pair of distinct oriented knot types, allowing its degree to depend on the pair?
\subsection{Sources}
\begin{description}
\item[\hypertarget{source:statement:1}{S1}] ULAM. AI and the UnsolvedMath Contributors, \emph{UnsolvedMath: A Curated Collection of Open Mathematics Problems}, catalogue entry. CC BY 4.0. \url{https://huggingface.co/datasets/ulamai/UnsolvedMath}.
\item[\hypertarget{source:statement:2}{S2}] T. Ohtsuki (editor), \emph{Problems on invariants of knots and 3-manifolds}, Geometry \& Topology Monographs 4 (2002), pp. 377--572 (published 2004), item 2.5 and its surrounding definitions and remarks. \url{https://arxiv.org/abs/math/0406190v1}.
\end{description}
\label{end:statement}
\end{document}
